\chapter{General Conclusion}

\vspace*{0.3cm}

This project focused on the optimization and compression perspective of Small
Language Models for Edge AI applications. The main objective was to prepare and
fine-tune \texttt{google/flan-t5-small} on an instruction-following dataset and
to compare different optimization algorithms according to both model quality and
resource efficiency.

\vspace*{0.15cm}

The first part of the work consisted of preparing the dataset
\texttt{yahma/alpaca-cleaned}. The original examples were cleaned, shuffled, and
converted into an instruction-tuning format using \texttt{source\_text} as the
model input and \texttt{target\_text} as the expected output. The processed data
was divided into training, validation, and test subsets, then saved in
\texttt{CSV} and \texttt{JSONL} formats.

\vspace*{0.15cm}

The second part of the project focused on the fine-tuning methodology. A training
pipeline was implemented to load the model, tokenize the dataset, select an
optimizer, run the training loop, evaluate the model, and save the resulting
metrics. The tested optimization methods included SGD, SGD with momentum, Adam,
AdamW with clipping and scheduling, L-BFGS, and Newton-CG.

\vspace*{0.15cm}

The experimental results showed a clear tradeoff between fine-tuning quality and
computational efficiency. Adam achieved the best training and validation losses,
which makes it the strongest optimizer in terms of model quality. However, it was
not the most efficient method in terms of latency and throughput. SGD and SGD with
momentum provided better efficiency, lower latency, and lower memory usage, making
them more suitable for constrained environments.

\vspace*{0.15cm}

The results also confirmed that L-BFGS and Newton-CG are less practical for this
type of transformer fine-tuning. Although these methods are important from a
theoretical optimization perspective, their high computational cost and memory
requirements make them difficult to use in an Edge AI context.

\vspace*{0.15cm}

Overall, the project demonstrates that optimizer selection must depend on the
target objective. If the priority is fine-tuning quality, Adam is the best
choice. If the priority is efficiency for Edge AI deployment, SGD with momentum or
AdamW with clipping and scheduling can offer a better compromise. This conclusion
is important because Edge AI systems require models that are not only accurate,
but also fast, compact, and resource-efficient.

\vspace*{0.15cm}

As future work, the fine-tuned models can be further optimized using compression
techniques such as quantization, pruning, or knowledge distillation. These methods
could reduce model size, memory usage, and inference latency while preserving an
acceptable level of performance. This would complete the transition from
optimizer comparison to practical deployment of Small Language Models on edge
devices.
